\begin{center}
\textbf{ABSTRACT}
\end{center}

$\!$\\

% O resumo em inglês deve ser organizado em apenas um parágrafo mesmo.

Traditional investment allocation methods often produce unstable portfolios because they rely solely on historical averages. This limitation highlights the importance of approaches that make the allocation process more consistent under financial market uncertainty. In this context, this study proposes an automated system to optimize a portfolio of up to 20 highly liquid stocks listed on the Brazilian stock exchange (B3), using the Black–Litterman model. The model combines return forecasts generated by machine learning (XGBoost) with AI-based sentiment analysis of financial headlines (GPT) to formulate investment views.

These views are integrated into market equilibrium through a dynamic confidence mechanism that adjusts the weight assigned to each view based on the model’s historical predictive performance and the volume of news available each day. The system operates end to end—from data collection to final portfolio allocation—allowing only long positions and imposing a maximum allocation limit per company.

Evaluated through simulations using data unseen during training and accounting for realistic transaction costs, the model was compared with a classical portfolio optimization approach (Markowitz) and the Selic benchmark interest rate. The results suggest that integrating artificial intelligence, automated text analysis, and statistical risk control offers a practical and promising approach to supporting investment decisions in the Brazilian stock market.

\vspace{1cm}

\hspace{-1.3cm}Keywords: Portfolio Optimization; Black-Litterman Model; Machine Learning; Sentiment Analysis; Stock Market.



%% pq a previsão melhora quando vamos pra semanas e mes ao inves de dias (quantidade de noticias + diversos outros fatores)